\section{Structural Normalization for the Next Companion Stage}
\label{sec:structural-normalization}

\begin{remark}[Version note]
Version \docversion\ introduces a structural normalization layer for the
Companion. The Millennium matrix is read primarily as the global matrix
realization of the Heisenberg compensator; the tool apparatus is split into
general Millennium-matrix tools and special Heisenberg-compensator tools; the
carrier/group/physical-readout distinction is made explicit; and the global
reader orientation is organized through a master chain from triolic source to
field-equation readout.
\end{remark}

\begin{remark}[Null axis versus light axis]
In structurally prior parts of the Companion we use the term \emph{null axis}
for the distinguished compensated axis selected by the MM/HC architecture.
When the same axis is read at the physical propagation level, it may be called
the \emph{light axis}. Thus ``light axis'' is not a second object, but the
physical readout of the structurally prior null axis.
\end{remark}

